% Part: first-order-logic
% Chapter: models-theories
% Section: expressing-relations

\documentclass[../../../include/open-logic-section]{subfiles}

\begin{document}

\olfileid{fol}{mat}{exr}

\olsection{بیان روابط در \article{structure}
  \printtoken{S}{structure}}

\begin{explain}
یکی از کاربردهای اصلی !!{formula}s بیان خواص و روابط در
!!a{structure}~$\Struct M$ بر حسب مفاهیم اولیهٔ زبان~$\Lang L$یِ
~$\Struct M$ است. منظورمان چنین است: !!{domain} ساختار~$\Struct M$ مجموعه‌ای
از اشیاست. !!{constant}s، !!{function}s و !!{predicate}s در~$\Struct M$
به‌ترتیب با اشیایی در~$\Domain M$، توابعی روی~$\Domain M$ و روابطی
روی~$\Domain M$ تعبیر می‌شوند. برای نمونه، اگر~$\Obj A^2_0$ در~$\Lang L$
باشد، آنگاه $\Struct M$ رابطه‌ای چون~$R = \Assign{{\Obj
  A^2_0}}{M}$ را به آن نسبت می‌دهد. در این صورت فرمول
$\Atom{\Obj A^2_0}{\Obj v_1, \Obj v_2}$ همان رابطه را به معنای زیر
\emph{بیان می‌کند}: اگر گمارش متغیر~$s$، متغیر~$\Obj v_1$ را به
$a \in \Domain{M}$ و~$\Obj v_2$ را به~$b \in \Domain M$ بفرستد، آنگاه
\[
Rab \text{\quad هرگاه و تنها هرگاه\quad} \Sat{M}{\Atom{\Obj A^2_0}{\Obj v_1, \Obj v_2}}[s].
\]
توجه کنید که در اینجا باید گمارش‌های متغیر را دخیل کنیم: نمی‌توانیم
صرفاً بگوییم «$Rab$ هرگاه و تنها هرگاه
$\Sat{M}{\Atom{\Obj A^2_0}{a, b}}$»، زیرا $a$ و~$b$ نمادهای زبان ما
نیستند، بلکه !!{element}s دامنهٔ~$\Domain{M}$ هستند.

ازآنجاکه فقط !!{formula}s اتمی نداریم، بلکه می‌توانیم آن‌ها را با
رابط‌های منطقی و سورها ترکیب کنیم، !!{formula}s پیچیده‌تر می‌توانند
روابط دیگری را تعریف کنند که مستقیماً در~$\Struct M$ گنجانده نشده‌اند.
می‌خواهیم بدانیم این کار چگونه انجام می‌شود و، به‌ویژه، کدام روابط را
می‌توانیم در !!a{structure} تعریف کنیم.
\end{explain}

\begin{defn}
فرض کنید $!A(\Obj v_1,\dots, \Obj v_n)$ !!a{formula} از~$\Lang L$ باشد
که در آن فقط $\Obj v_1$،\dots، $\Obj v_n$ آزاد رخ می‌دهند و فرض کنید
$\Struct M$ !!a{structure} برای~$\Lang L$ باشد. فرمول
$!A(\Obj v_1,\dots, \Obj v_n)$ رابطهٔ
$R \subseteq \Domain M^n$ را \emph{بیان می‌کند} هرگاه و تنها هرگاه
\[
Ra_1\dots a_n \text{\quad هرگاه و تنها هرگاه\quad} \Sat{M}{\Atom{!A}{\Obj
    v_1,\dots, \Obj v_n}}[s]
\]
برای هر گمارش متغیر~$s$ که $s(\Obj v_i) = a_i$ ($i = 1,
\dots, n$).
\end{defn}

\begin{ex}
در مدل استاندارد حساب~$\Struct N$، !!{formula}~$\Obj
v_1 < \Obj v_2 \lor \eq[\Obj v_1][\Obj v_2]$ رابطهٔ~$\le$ را روی~$\Nat$
بیان می‌کند. !!{formula}~$\eq[\Obj v_2][\Obj v_1']$ رابطهٔ جانشینی، یعنی
رابطهٔ~$R \subseteq \Nat^2$ را بیان می‌کند که در آن $Rnm$ برقرار است اگر
$m$ جانشین~$n$ باشد. فرمول~$\eq[\Obj v_1][\Obj v_2']$ رابطهٔ پیشینی را
بیان می‌کند. دو !!{formula}~$\lexists[\Obj v_3][(\eq/[\Obj v_3][\Obj 0] \land
  \eq[\Obj v_2][(\Obj v_1 + \Obj v_3)])]$ و~$\lexists[\Obj
  v_3][\eq[(\Obj v_1 + {\Obj v_3}')][\Obj v_2]]$ هر دو رابطهٔ~$\Obj <$ را
بیان می‌کنند. معنای این امر آن است که نماد محمول~$<$ در واقع
در زبان حساب زائد است؛ می‌توان آن را تعریف کرد.
\end{ex}

\begin{explain}
این ایده فقط در !!{structure}s خاص جالب نیست، بلکه هرگاه از زبانی برای
توصیف یک یا چند مدل مورد نظر استفاده می‌کنیم، یعنی هنگامی که نظریه‌ها
را بررسی می‌کنیم، به‌طور کلی اهمیت دارد. این نظریه‌ها اغلب فقط چند
!!{predicate} را به‌عنوان نمادهای پایه در بر می‌گیرند، اما در حوزه‌ای
که برای توصیف آن به کار می‌روند، بسیاری از روابط دیگر نیز نقشی مهم
دارند. اگر بتوان این روابط دیگر را به‌طور نظام‌مند با روابطی بیان کرد
که !!{predicate}s پایهٔ زبان را تعبیر می‌کنند، می‌گوییم می‌توانیم آن‌ها
را در زبان \emph{تعریف} کنیم.
\end{explain}

\begin{prob}
!!{formula}sی در~$\Lang L_A$ بیابید که روابط زیر را تعریف کنند:
\begin{enumerate}
\item $n$ میان $i$ و~$j$ است؛
\item $n$، عدد~$m$ را بی‌باقی‌مانده تقسیم می‌کند (یعنی $m$ مضربی از~$n$ است)؛
\item $n$ عددی اول است (یعنی هیچ عددی جز~$1$ و~$n$، عدد~$n$ را
  بی‌باقی‌مانده تقسیم نمی‌کند).
\end{enumerate}
\end{prob}

\begin{prob}
فرض کنید فرمول~$!A(\Obj v_1, \Obj v_2)$ رابطهٔ
$R \subseteq \Domain M^2$ را در !!a{structure}~$\Struct M$ بیان کند.
فرمول‌هایی بیابید که روابط زیر را بیان کنند:
\begin{enumerate}
\item وارون~$R^{-1}$ از~$R$؛
\item حاصل‌ضرب نسبی~$R \mid R$؛
\end{enumerate}
آیا می‌توانید راهی برای بیان~$R^+$، یعنی بستار تراگذری~$R$، بیابید؟
\end{prob}

\begin{prob}
فرض کنید $\Lang{L}$ زبانی باشد که فقط شامل نماد محمول
دوموضعی~$<$ است (و هیچ !!{constant}، !!{function} یا !!{predicate} دیگری
ندارد---البته به‌جز~$\eq$). فرض کنید $\Struct{N}$ ساختاری باشد چنان‌که
$\Domain{N} = \Nat$ و~$\Assign{<}{N} = \Setabs{\tuple{n,m}}{n <
  m}$. موارد زیر را اثبات کنید:
\begin{enumerate}
\item $\{ 0 \}$ در~$\Struct{N}$ تعریف‌پذیر است؛
\item $\{ 1 \}$ در~$\Struct{N}$ تعریف‌پذیر است؛
\item $\{ 2 \}$ در~$\Struct{N}$ تعریف‌پذیر است؛
\item برای هر $n \in \Nat$، مجموعهٔ~$\{ n \}$ در~$\Struct{N}$
  تعریف‌پذیر است؛
\item هر زیرمجموعهٔ متناهیِ~$\Domain{N}$ در~$\Struct{N}$ تعریف‌پذیر است؛
\item هر زیرمجموعهٔ هم‌متناهیِ~$\Domain{N}$ در~$\Struct{N}$ تعریف‌پذیر
  است (که در آن $X \subseteq \Nat$ هم‌متناهی است هرگاه و تنها هرگاه
  $\Nat \setminus X$ متناهی باشد).
\end{enumerate}
\end{prob}


\end{document}
